\chapter{Iloczyny nieskończone}
\label{ma1:ch:infinite-products}

W definicji szeregu przechodziliśmy do granicy sum częściowych. Analogicznie można badać granice iloczynów częściowych. Istnieje jednak ważna różnica: dodanie skończenie wielu składników nie może zmienić zbieżności szeregu, natomiast jeden czynnik równy zeru sprawia, że wszystkie dalsze iloczyny częściowe są zerowe. Odróżnimy zatem zbieżność do niezerowej granicy od zanikania iloczynów częściowych. Dowody przeprowadzimy bez odwoływania się do logarytmu.

\section{Definicja i własności elementarne}
\begin{definicja}
\label{ma1:def:infinite-product}
Niech $K=\R$ lub $K=\C$, a $(u_n)_{n\ge1}$ będzie ciągiem w $K$. Iloczyn nieskończony $\prod_{n=1}^\infty u_n$ rozpatrujemy wraz z ciągiem iloczynów częściowych
\[
 P_0=1,\qquad P_N=\prod_{n=1}^N u_n\quad(N\in\N_+).
\]
Jeżeli wszystkie $u_n\ne0$, to iloczyn nazywamy zbieżnym, gdy istnieje granica $P_N\to P\in K\setminus\{0\}$. Wówczas piszemy $\prod_{n=1}^\infty u_n=P$. W przeciwnym razie, przy niezerowych czynnikach, iloczyn nazywamy rozbieżnym, również gdy $P_N\to0$.
\end{definicja}
Przyjmujemy w tym rozdziale konwencję zbieżności do niezerowej granicy. Gdy występuje czynnik zerowy, wszystkie iloczyny częściowe od jego indeksu są zerowe; tę sytuację rozpatrujemy osobno. Można również stosować konwencję dopuszczającą skończenie wiele czynników zerowych i wymagającą zbieżności niezerowego ogona. W każdym przypadku trzeba odróżnić zwykłą granicę ciągu $(P_N)$ od zbieżności iloczynu w powyższym sensie.

\begin{twierdzenie}
\label{ma1:thm:product-elementary}
Niech $(u_n),(v_n)$ będą ciągami niezerowych liczb z $K=\R$ lub $K=\C$.
\begin{enumerate}[label=(\roman*)]
\item Dla dowolnego $m\in\N_+$ iloczyn $\prod_{n\ge1}u_n$ jest zbieżny wtedy i tylko wtedy, gdy $\prod_{n\ge m}u_n$ jest zbieżny. W przypadku zbieżności
$\prod_{n\ge1}u_n=(\prod_{n=1}^{m-1}u_n)(\prod_{n\ge m}u_n)$.
\item Jeżeli $\prod u_n=U$ i $\prod v_n=V$ są zbieżne, to $\prod(u_nv_n)=UV$ oraz $\prod u_n^{-1}=U^{-1}$.
\item Zbieżność $\prod u_n$ implikuje $u_n\to1$.
\end{enumerate}
\end{twierdzenie}
\begin{proof}
Dla (i) iloczyn częściowy od $m$ do $N$ wynosi $P_N/P_{m-1}$, gdzie $P_{m-1}\ne0$. Mnożenie przez niezerową stałą i dzielenie przez nią zachowuje istnienie i niezerowość granicy. Dla (ii) oznaczmy $Q_N=\prod_{n=1}^N v_n$. Iloczyny częściowe szeregów czynników $u_nv_n$ oraz $u_n^{-1}$ są odpowiednio $P_NQ_N$ i $1/P_N$. Rachunek granic daje granice $UV\ne0$ i $1/U\ne0$. Dla (iii) mamy $u_n=P_n/P_{n-1}$. Oba ciągi w ilorazie dążą do tego samego niezerowego $P$, zatem $u_n\to P/P=1$.
\end{proof}
\begin{przyklad}
\label{ma1:ex:telescoping-products}
Dla $n\ge2$ rozważmy trzy ciągi czynników. Z bezpośredniego skracania iloczynów skończonych dla $N\ge2$ otrzymujemy
\[
\begin{aligned}
 \prod_{n=2}^N(1-1/n)&=1/N,\\
 \prod_{n=2}^N(1+1/n)&=(N+1)/2,\\
 \prod_{n=2}^N(1-1/n^2)&=
 \prod_{n=2}^N\frac{n-1}{n}\frac{n+1}{n}
 =\frac{N+1}{2N}.
\end{aligned}
\]
W każdym przypadku czynniki dążą do $1$. W pierwszym przypadku iloczyny częściowe dążą do zera, w drugim do $+\infty$, a tylko w trzecim iloczyn jest zbieżny w sensie definicji \ref{ma1:def:infinite-product}, do $1/2$. Warunek konieczny zbieżności nie jest wystarczający.
\end{przyklad}

\section{Kryterium Cauchy'ego dla iloczynów}
\begin{twierdzenie}
\label{ma1:thm:product-cauchy}
Niech $(u_n)_{n\ge1}$ będzie ciągiem niezerowych liczb rzeczywistych lub zespolonych. Iloczyn $\prod_{n\ge1}u_n$ jest zbieżny wtedy i tylko wtedy, gdy
\[
 \forall\varepsilon\in\R_{>0}\ \exists N\in\N_+\
 \forall q,p\in\N_+\quad
 \left(q\ge p\ge N\ \Rightarrow\
 \left|\prod_{n=p}^{q}u_n-1\right|<\varepsilon\right).
\]
\end{twierdzenie}
\begin{proof}
Jeżeli $P_n\to P\ne0$, to od pewnego indeksu $|P_n-P|<|P|/2$, więc $|P_n|\ge|P|/2$. Dla $\varepsilon>0$ wybieramy $M$ tak, aby dla wszystkich $n\ge M$ zachodziły zarówno ta nierówność, jak i $|P_n-P|<\varepsilon|P|/4$. Gdy $q\ge p\ge M+1$, mamy
\[
 \left|\prod_{n=p}^q u_n-1\right|
 =\frac{|P_q-P_{p-1}|}{|P_{p-1}|}
 \le\frac{|P_q-P|+|P_{p-1}-P|}{|P|/2}
 <\varepsilon.
\]

Odwrotnie, stosujemy warunek z tezy najpierw dla $\varepsilon=1/2$. Istnieje $N_0\ge1$ takie, że dla każdego $q\ge N_0$ mamy $|\prod_{n=N_0}^q u_n-1|<1/2$. Z nierówności trójkąta i nierówności odwrotnej wynika
\[
 \frac12|P_{N_0-1}|<|P_q|<\frac32|P_{N_0-1}|.
\]
Ciąg iloczynów częściowych jest więc ograniczony, również po uwzględnieniu skończonego początku; niech $B\ge1$ będzie ograniczeniem modułów wszystkich $P_n$ dla $n\ge0$.

Dla dowolnego $\varepsilon>0$ wybieramy $N$ z warunku tezy dla dokładności $\varepsilon/B$. Dla $q>p\ge N$ otrzymujemy
\[
 |P_q-P_p|=|P_p|\left|\prod_{n=p+1}^q u_n-1\right|
 <B\frac{\varepsilon}{B}=\varepsilon.
\]
Przypadek $p=q$ jest trywialny, a dla $q<p$ zamieniamy indeksy. Z zupełności $\R$ i $\C$ ciąg $(P_n)$ ma granicę $P$. Dolne oszacowanie $|P_q|>|P_{N_0-1}|/2$ dla $q\ge N_0$ i zbieżność modułów dają $|P|\ge |P_{N_0-1}|/2>0$. Zatem iloczyn jest zbieżny.
\end{proof}

\section{Oszacowania skończonych iloczynów}
\begin{lemat}
\label{ma1:lem:product-finite-estimate}
Niech $m\in\N_+$ i niech $z_1,\ldots,z_m\in K$, gdzie $K=\R$ lub $K=\C$. Oznaczmy $s=\sum_{j=1}^m|z_j|$. Wówczas
\[
 \left|\prod_{j=1}^m(1+z_j)-1\right|
 \le \prod_{j=1}^m(1+|z_j|)-1.
\]
Jeżeli $s<1$, to
\[
 \prod_{j=1}^m(1+|z_j|)\le\frac1{1-s},
 \qquad
 \left|\prod_{j=1}^m(1+z_j)-1\right|
 \le\frac{s}{1-s}.
\]
\end{lemat}
\begin{proof}
Rozwijając iloczyn $\prod(1+z_j)$ za pomocą rozdzielności, otrzymujemy sumę po wszystkich podzbiorach zbioru $\{1,\ldots,m\}$. Składnik odpowiadający zbiorowi pustemu wynosi $1$. Po odjęciu go i zastosowaniu nierówności trójkąta do pozostałych składników dostajemy pierwsze oszacowanie, gdyż suma ich modułów to dokładnie $\prod(1+|z_j|)-1$.

Niech $t_j=|z_j|\ge0$ oraz $s_k=\sum_{j=1}^k t_j$ dla $0\le k\le m$, gdzie $s_0=0$. Jeśli $s_m<1$, to $0\le s_k<1$ dla każdego $k$. Pokażemy indukcyjnie, że $\prod_{j=1}^k(1+t_j)\le1/(1-s_k)$. Dla $k=0$ zachodzi równość. W kroku indukcyjnym korzystamy z
\[
 (1-s_k)(1+t_k)
 =(1-s_{k-1}-t_k)(1+t_k)
 =1-s_{k-1}-s_{k-1}t_k-t_k^2
 \le1-s_{k-1}.
\]
Ponieważ oba mianowniki są dodatnie, wynika stąd $(1+t_k)/(1-s_{k-1})\le1/(1-s_k)$. To daje krok indukcyjny. Ostatnie oszacowanie wynika z pierwszego oraz $1/(1-s)-1=s/(1-s)$.
\end{proof}

\begin{lemat}
\label{ma1:lem:positive-product-bound}
Niech $t_1,\ldots,t_m$ będą nieujemnymi liczbami rzeczywistymi. Wówczas
$\prod_{j=1}^m(1+t_j)\ge1+\sum_{j=1}^m t_j$.
Jeżeli ponadto $t_j\le1$ dla każdego $j$, to
$\prod_{j=1}^m(1-t_j)\ge1-\sum_{j=1}^m t_j$.
\end{lemat}
\begin{proof}
Pierwsza nierówność wynika z rozwinięcia iloczynu: składniki stopnia co najmniej dwa są nieujemne. Dla drugiej przeprowadzamy indukcję. Gdy $m=1$, zachodzi równość. Z założenia indukcyjnego, po pomnożeniu przez nieujemne $1-t_{m+1}$, mamy
\[
 \prod_{j=1}^{m+1}(1-t_j)
 \ge (1-\sum_{j=1}^m t_j)(1-t_{m+1})
 =1-\sum_{j=1}^{m+1}t_j+t_{m+1}\sum_{j=1}^m t_j
 \ge1-\sum_{j=1}^{m+1}t_j.
\]
\end{proof}

\section{Zbieżność bezwzględna}
\begin{definicja}
\label{ma1:def:absolute-product}
Niech $(a_n)_{n\ge1}$ będzie ciągiem liczb rzeczywistych lub zespolonych takim, że $1+a_n\ne0$ dla każdego $n$. Iloczyn $\prod_{n\ge1}(1+a_n)$ nazywamy bezwzględnie zbieżnym, jeżeli szereg $\sum_{n\ge1}|a_n|$ jest zbieżny.
\end{definicja}
\begin{twierdzenie}
\label{ma1:thm:absolute-product}
Niech $(a_n)$ spełnia założenia definicji \ref{ma1:def:absolute-product}. Jeśli $\sum|a_n|$ jest zbieżny, to $\prod(1+a_n)$ jest zbieżny do niezerowej granicy $P$. Ponadto, oznaczając $P_N=\prod_{n=1}^N(1+a_n)$ i $P_0=1$, dla każdego całkowitego $N\ge0$ takiego, że $s_N=\sum_{n>N}|a_n|<1$, otrzymujemy
\[
 \left|\frac{P}{P_N}-1\right|\le\frac{s_N}{1-s_N}.
\]
\end{twierdzenie}
\begin{proof}
Ustalmy $\varepsilon>0$ i wybierzmy rzeczywiste $\delta$ z $0<\delta<\min(1/2,\varepsilon/2)$. Ponieważ szereg modułów jest zbieżny, istnieje $N_0$ z $\sum_{n\ge N_0}|a_n|<\delta$. Dla $q\ge p\ge N_0$ mamy $s=\sum_{n=p}^q|a_n|<\delta$. Z lematu \ref{ma1:lem:product-finite-estimate} wynika
\[
 \left|\prod_{n=p}^q(1+a_n)-1\right|
 \le\frac{s}{1-s}\le2s<\varepsilon.
\]
Kryterium \ref{ma1:thm:product-cauchy} daje zbieżność iloczynu do niezerowego $P$.

Dla ustalonego $N$ z $s_N<1$ oraz $q>N$ mamy $\sum_{n=N+1}^q|a_n|\le s_N$, a funkcja $s\mapsto s/(1-s)$ jest rosnąca na $[0,1)$, gdyż dla $0\le s\le t<1$ nierówność $s/(1-s)\le t/(1-t)$ po pomnożeniu przez dodatnie mianowniki jest równoważna $s\le t$. Z lematu otrzymujemy $|P_q/P_N-1|\le s_N/(1-s_N)$. Przechodząc do granicy po $q$, dostajemy tezę.
\end{proof}
\begin{przyklad}
Dla $a_n=2^{-n}$ mamy $s_N=2^{-N}$, więc dla $N\ge1$ iloczyn $\prod_{n=1}^\infty(1+2^{-n})$ spełnia
$|P/P_N-1|\le1/(2^N-1)$. Dla $N=10$ jest to $1/1023<1/1000$. Możemy więc oszacować błąd względny względem obliczonego iloczynu częściowego bez znajomości dokładnej wartości $P$.
\end{przyklad}
